%% LyX 2.4.4 created this file.  For more info, see https://www.lyx.org/.
%% Do not edit unless you really know what you are doing.
\documentclass{article}
\usepackage{amsmath}
\usepackage{amssymb}
\usepackage{fontspec}
\setlength{\parindent}{0bp}
\usepackage{color}
\usepackage{pifont}
\usepackage{float}
\usepackage{units}
\usepackage{cancel}
\usepackage{geometry}
\geometry{verbose,tmargin=2.7cm,bmargin=2cm,lmargin=2cm,rmargin=2cm,headheight=2cm,headsep=0cm,footskip=0cm}
\usepackage{setspace}
\PassOptionsToPackage{normalem}{ulem}
\usepackage{ulem}
\usepackage{microtype}
\onehalfspacing
\usepackage[pdfusetitle,
 bookmarks=true,bookmarksnumbered=true,bookmarksopen=true,bookmarksopenlevel=2,
 breaklinks=true,pdfborder={0 0 0},pdfborderstyle={},backref=section,colorlinks=true]
 {hyperref}
\hypersetup{
 linkcolor=black, citecolor=black}

\makeatletter

%%%%%%%%%%%%%%%%%%%%%%%%%%%%%% LyX specific LaTeX commands.
\providecommand\textquotedblplain{%
  \bgroup\addfontfeatures{Mapping=}\char34\egroup}

\@ifundefined{date}{}{\date{}}
%%%%%%%%%%%%%%%%%%%%%%%%%%%%%% User specified LaTeX commands.
% Engine guard: use XeLaTeX/LuaLaTeX
\usepackage{iftex}
\ifPDFTeX
  \PackageError{template}{Please compile with XeLaTeX or LuaLaTeX}{}
\fi

\begin{document}

\section{שבוע 2 - הרצאה 3 28 אוקטובר 2025}

\subsection*{תקציר השבוע הקודם}
\begin{itemize}
\item פתרון התכונות המאקרוסקופיות של מערכת מרובת חלקיקים - שיטת ההסתברויות.
\item חוק המספרים הגדולים - ההתפלגות בפועל שווה להתפלגות החזויה.
\item משפט הגבול המרכזי - כאשר מספר המשתנים האקראיים מאוד גדול ההסתברות
של הסכום מתרכזת סביב הערך הסביר ביותר.
\end{itemize}
תאור מערכות פיזיקליות:

מיקרו-מצבים - תאור פרטני של כל דרגות החופש.

מאקרו-מצב - אפיון המערכת לפי תכונות מאקרוסקופיות.

כפליות - מספר המיקרו-מצבים המתאימים למאקרו-מצב.

\textbf{\textcolor{red}{הנחה: ההסתברות של כל המיקרו-מצבים המותרים
שווה.}}

$\Leftarrow$ ההסתברות של מאקרו-מצב פרופרוציונית לכפליות שלו.
מסקנה : המערכת נוטה להגיע למאקרו-מצב הסביר ביותר (ושכניו הקרובים אליו)
ואז להשאר בו - מצב ש\textquotedbl מ.
אנטרופיה

\[
S=k_{B}\ln\left(\Omega\right)
\]

\[
\Delta S\geqslant0
\]


\subsubsection*{מוצק אינשטיין:}

$N+1$ חלקיקים הניתנים לאבחנה ו-$q$ מנות אנרגיה.
\[
\Omega\left(N+1,q\right)=\frac{\left(N+q\right)!}{q!N!}
\]
\[
\Omega\left(N+1,q_{A},N+1,q_{B}\right)=\Omega\left(N+1,q_{A}\right)\Omega\left(N+1,q_{B}\right)
\]

\[
\ln\left(\Omega\left(N+1,q_{A},N+1,q_{B}\right)\right)=\ln\left(\Omega\left(N+1,\frac{q}{2},N+1,\frac{q}{2}\right)e^{-\left(\frac{x}{\nicefrac{q}{2}}\right)^{2}N}\right)
\]
 

המצב הסביר ביותר הוא $\frac{q}{2}=q_{A}=q_{B}$

הסיכוי לסטות מהמצב הסביר ביותר ב-$x$

\[
P\left(x\neq0\right)=P\left(x=0\right)\cdot\exp\left(-\left(\frac{x}{\nicefrac{q}{2}}\right)^{2}\right)N
\]


\end{document}
